\documentclass[10pt]{article}
\usepackage{microtype}
\usepackage{graphicx}
\usepackage{subcaption}
\usepackage{booktabs}
\usepackage{xurl}
\usepackage{hyperref}
\usepackage[T1]{fontenc}
\usepackage{tmlr}
\usepackage{amsmath,amssymb,amsfonts,mathtools}
\usepackage{amsthm}
\usepackage{placeins}
\usepackage[nameinlink,capitalize,noabbrev]{cleveref}
\newcommand{\R}{\mathbb{R}}
\newcommand{\Id}{\mathbf{I}}
\newcommand{\norm}[1]{\left\lVert #1\right\rVert}
\newcommand{\Eeq}{E_{\mathrm{eq}}}
\newcommand{\Ehom}{E_{\mathrm{hom,data}}}
\newcommand{\Ehomop}{E_{\mathrm{hom,op}}}
\newcommand{\Erel}{E_{\mathrm{rel,data}}}
\newcommand{\Erelop}{E_{\mathrm{rel,op}}}
\newcommand{\Ecomp}{E_{\mathrm{comp,data}}}
\newcommand{\Ecompop}{E_{\mathrm{comp,op}}}
\newtheorem{proposition}{Proposition}
\newtheorem{corollary}[proposition]{Corollary}
\newtheorem{spectraltheorem}[proposition]{Theorem}
\newtheorem{spectrallemma}[proposition]{Lemma}
\crefname{spectraltheorem}{Theorem}{Theorems}
\crefname{spectrallemma}{Lemma}{Lemmas}
\title{Beyond Shared Geometry: Predictive Recovery, Exact Actions, and Semantic Specificity in Language-Model Representations}
\author{Anonymous authors}
\input{generated/bounded_macros.tex}
\begin{document}
\maketitle

\begin{abstract}
What can recovered transformations reveal about language-model representations beyond predictive accuracy? Given a prescribed finite-group action on input states, we fit maps that predict hidden representations of transformed states from their source representations in the same context. Each operator family is shared across contexts within a fixed model-layer coordinate space. Our architecture-independent theory distinguishes optimal unconstrained linear prediction from minimum risk over fixed-dimensional exact orthogonal actions under a common second-moment-whitened metric. An analytical example gives two observation systems with identical optimal regression risk and eigenvalue spectra of their orbit Gram matrices, but different exact-action risks. A spectral budget respecting real irreducible representation types gives lower and upper risk bounds differing by at most a factor of two. We study weekday cycles ($C_7$) and square-corner actions ($D_4$) in six frozen, pretrained decoder-only autoregressive Transformers from the Llama and Qwen families. Linear recovery, label-informed category controls and compact nonlinear predictors test competing predictive explanations; the latter are not established nonlinear group actions. Finite-bank lower bounds on the exact orthogonal-action premium over independent linear prediction are numerically positive in all \BoundedCellCount{} conditions. Frozen witnesses have higher held-out loss than independent least squares in all 48 conditions, while exhaustive permutation-image comparisons do not uniformly favor the specified semantic image. On new authored contexts under matched validation selection, difference fitting lowers mean weekday loss but raises mean square-corner loss relative to coordinate fitting. These results distinguish predictive recovery, exact-action compatibility and semantic interpretation without identifying a uniquely supported semantic action or native computation.
\end{abstract}

\section{Introduction}
Which additional interpretation is justified when a fitted transformation predicts well, approximately obeys group laws, and outperforms disrupted correspondences? These tests may appear to reinforce one another, yet they address different questions. A useful predictor is not by itself an account of a model's computation; consistency can follow partly from accurate prediction; disrupting a registration does not test whether the specified semantic correspondence is uniquely favored. We study these distinctions through controlled finite actions in contextual language-model representations.

Transformation fitting is established: \citet{lenc2015understanding} study paired maps and network stitching, and \citet{sun2026transformationlaws} develops compatible transformation families, realizability and transport. Equivariance by Contrast identifies jointly learned coordinates and actions under explicit assumptions \citep{schmidt2025equivariance}. Calendar mechanisms are studied by \citet{engels2025notall}; group-operation mechanisms and competing explanations by \citet{nanda2023progress,chughtai2023toy,wu2025unified}. Implementation-family equivalence requires abstraction assumptions \citep{sun2026tracking}. We test which identification claims finite-action controls support; \cref{app:related-work} gives the comparisons.

The protocol fits a linear operator family within each fixed model, layer and training-derived coordinate space to seven weekday states and four corner states. Each operator predicts the representation of its transformed state from the source representation in the same context; the same family is then used across held-out templates without context-specific realignment. Matrices are fitted separately across models and layers. States are visible in the input, and the language models remain unchanged. Independent fitting removes group-law penalties from the objective; generator-only fitting withholds action elements. Their predictive success motivates comparisons with category predictors, a coherent alternative registration, compact nonlinear predictors and context-difference fits.

\input{generated/main_findings.tex}

The paper makes three connected contributions. First, it gives a computable bound on the best risk over the fixed-dimensional exact orthogonal-action class and separates that target from ordinary prediction error. Second, it connects this bound to measured representations through 48 spectral analyses and constructed witnesses, then compares all distinct small-group permutation images. Third, it tests whether an earlier context-transport advantage survives a new authored language construction under matched validation selection. The result is a specific empirical distinction: exact-action compatibility is widespread in this panel, semantic ranking is condition-dependent, and the preferred transport objective changes across domains. Standard finite-orbit and error-propagation identities serve as supporting interpretation rather than additional new theory.

The practical question is consequently not whether every recovery score is positive. It is whether a specific comparison changes the explanation supported by the data. Low prediction error establishes a useful transformation; a small exact-class risk establishes compatibility with one shared action under a specified metric; preference over coherent alternatives and causal use require different observations. This distinction yields a focused way to read the measurements and to decide which additional evaluation is informative.

\section{Protocol and Evidence Scope}
\label{sec:setup}
Let $G$ act on states $i$, and let $t$ index contextual templates. For each model, layer, and split, training activations determine mean $\mu$ and orthonormal PCA columns $C\in\R^{d\times q}$. We use
\begin{equation}
 z_{i,t}=(h_{i,t}-\mu)C,\qquad z_{i,t}R_g\approx z_{g(i),t}.
\end{equation}
For row vectors, applying $g_2$ then $g_1$ gives
\begin{equation}
 R_{g_2}R_{g_1}\approx R_{g_1\circ g_2}.
 \label{eq:convention}
\end{equation}
The primary dimension is $q=16$. The lexical-rank comparison (AFS-D in the source records) uses matched $q=6$ for seven weekdays and $q=3$ for four corners.

\paragraph{Recovery target.}
The labeled input action is given, rather than discovered. We distinguish the risk of a fitted predictor, the optimum over a declared action class, and identification of action parameters. Exact realizability determines a linear action only on the supported span; unobserved extensions remain free (\cref{app:contextual-realization}). One shared map in fixed coordinates is a stronger requirement than separate context-wise realizations. Even exact shared prediction can arise from a stable state code and leaves the generating mechanism unidentified.

\paragraph{Estimation.}
Ridge fits each element independently, without an intercept, using regularization $\alpha=0.01$. Joint fitting starts from ridge: the data term has coefficient one, and the homomorphism, relation, identity and orthogonality penalties have coefficients $1,1,1,0.05$, respectively. AdamW runs for 1,200 steps, learning rate $0.002$, weight decay $10^{-5}$, and technical seed 42. All transformations use the same fixed training-only PCA. Full objectives appear in \cref{app:methods}.

\paragraph{Evaluation.}
For paired row matrices define
\begin{equation}
 D(A,B)=\frac{1}{m}\sum_{j=1}^m
 \frac{\norm{A_j-B_j}_2^2}{\norm{B_j}_2^2+10^{-12}}.
 \label{eq:D}
\end{equation}
$\Eeq$ averages $D(ZR_g,Z_g)$ over nonidentity actions on held-out templates. $\Ehom$ averages $D(ZR_{g_2}R_{g_1},ZR_{g_1\circ g_2})$ over element pairs. $\Ecomp$ uses the same expression excluding pairs containing the identity; it is a restricted consistency summary, not an independent unseen-element test. AFS-A instead compares generator products with \emph{actual} held-out transformed activations. Relation and operator-space definitions are explicit in \cref{app:metrics}. Data defects are squared relative errors, whereas the reported homomorphism, composition, and relation operator diagnostics use unsquared relative Frobenius norms. Their raw magnitudes cannot be directly compared.

\paragraph{Panel and partitions.}
\label{sec:design}
We use six frozen, pretrained base decoder-only autoregressive Transformer models---Llama-3.2-1B/3B, Llama-3.1-8B, and Qwen2.5-1.5B/3B/7B---at normalized depth slots $0.2,0.4,0.6,0.8$. The $C_7$ action cycles Monday through Sunday. The $D_4$ action rotates and reflects the labels NW, NE, SE, SW. Each domain contains 40 short templates, such as ``The current day is Monday.'' and ``The marked corner is NW.'' The state is visible in the input; no held-out vocabulary, unseen semantic domain, or abstract reasoning task is tested.

For each of five split seeds $(101,303,909,2701,4243)$, 28 templates train PCA and operators, six are diagnostic validation templates, and six are test templates. Complete state orbits stay within their template split. The same 40 templates recur across splits; the seeds represent repeated partitions, not five independently sampled datasets. Activations are mean hidden states over the unique lexical target-token span at transformer block outputs. Exact model revisions, tokenization, layers, prompts, and sensitivity subsets are in \cref{app:methods}.

\paragraph{Evidence panels.}
Study~1 is the structural recovery panel, called FULL and AFS in the source records. It compares ridge, joint, identity and random controls, then evaluates generator products, correspondence disruption and rank-matched lexical codes. Study~2 is the predictor and context-transport panel, using separate training, validation and test banks from a fixed grammar construction. Its C7/D4 test bank has 64 context families with two paired variants per family, yielding 128 complete state blocks per model/readout/domain. It does not reuse the five-partition design as five independent replications. The completed follow-up uses these retained coordinates for spectral, witness and image-subgroup comparisons, and adds a new 64-family authored context bank for transport evaluation. All 48 declared conditions and all 36 mechanism settings are reported; no historical calibration or native failure is superseded. The coordinate-only extensions use training-energy normalizations and explicit storage budgets. Detailed source mappings and preserved stopping decisions appear in the appendices.

\begin{table}[t]
\centering\small
\begin{tabular}{@{}p{.13\textwidth}p{.39\textwidth}p{.40\textwidth}@{}}
\toprule
Symbol & Definition or normalization & Target and scope \\
\midrule
$D(A,B)$ & Mean of squared error divided by each target row's squared norm & Study~1 paired prediction and consistency; the target in the denominator depends on the endpoint. \\
$E_z$ & $\mathbb E_{\rm train}\|z-\bar z_{\rm train}\|_2^2$ & Study~2 coordinate errors share this training-only denominator. \\
$E_u$ & $\mathbb E_{\rm train}\|u\|_2^2$, $u=z_{p_0}-z_{p_1}$ & Study~2 paired context-difference errors; no extra recentering. \\
$\eta,L^\star$ & Squared error after second-moment whitening, divided by $r$, with a full uniform group average & Same distribution and metric; independent least squares versus exact orthogonal actions. \\
$L(I),L^\star-\eta$ & Identity-action risk and exact-class premium in the same whitened metric & $I$ means $\rho_g=\Id_r$ for every $g$; the premium concerns exact laws and orthogonality jointly. \\
$\gamma_r$ & Type-constrained invariant-subspace energy budget divided by $r$ & A finite-bank quantity or its specified population counterpart; not a historical error score. \\
\bottomrule
\end{tabular}
\caption{Metrics and targets. Equal-looking error values under different rows of this table are not directly comparable.}
\label{tab:metric-contract}
\end{table}

\paragraph{Statistical scope.}
All completed empirical comparisons are descriptive analyses of the recorded fixed panels. Computational replay re-executes inspected designs and supplies no independent confirmation. We retain the recorded means and bootstrap percentile intervals. Historical FULL uses 5,000 draws, base seed 314159 with analysis-specific offsets, and model, depth, split, and template levels, adding random-control seed where relevant. AFS-D keeps depth fixed for its matched contrasts; the separate global-additive experiment resamples split aggregates. The six models are deliberately selected, layers are repeated measurements, and template partitions overlap. These intervals describe resampling sensitivity under the recorded scheme, rather than calibrated population coverage. Numerous secondary intervals are not globally multiplicity-adjusted. Counts such as 9/24 summarize fixed model/readout conditions, not independent Bernoulli trials. We make no new significance or equivalence decisions from these intervals; \cref{app:inference} gives their scope. The new-context contrast uses its separately frozen paired-family scheme, detailed in \cref{app:pending-protocols}; it does not replace the older panel retrospectively.

\section{Prediction Error and Exact-Action Risk}
\label{sec:boundaries}
For a fixed distribution and loss, a predictor's risk upper-bounds the optimum over any function class containing it. It does not determine the optimum over a different class constrained to obey group laws. We make this distinction precise, then relate it to finite-state recovery controls.

\subsection{Same summaries, different exact-action risk}
Consider complete labeled states $s\in S$ in context $c$, with fixed-coordinate observations $z_{c,s}\in\R^{1\times r}$. For a fixed context distribution $P_C$, average uniformly over states and all group elements, including identity and fixed points. For positive-definite second moment $\Sigma=\mathbb E[z^\top z]$, set $x=z\Sigma^{-1/2}$ and $x_g(c,s)=x(c,g\cdot s)$. Let $\eta$ be the minimum normalized squared prediction error for independently chosen matrices $A_g$, and let
\[
 L^\star=\inf_{\rho\in\mathcal R_r(G)}
 \frac1{r|G|}\sum_{g\in G}\mathbb E\norm{x\rho_g-x_g}_2^2,
\]
where $\mathcal R_r(G)$ contains exact row actions with $\rho_g\in O(r)$ and $\rho_g\rho_h=\rho_{hg}$. Both risks use this same whitened metric.

\paragraph{An analytical separation.}
Fix $G=S=C_2\times C_2\times C_7$ acting by addition, uniform states $(a,b,k)$, one context and $r=1$. The observations
\[
 x_A=\frac{(-1)^a+(-1)^b}{\sqrt2},\qquad
 x_B=\sqrt2\cos(2\pi k/7)
\]
have unit second moment, the same optimal regression risk $\eta=1/2$, and the same untyped nonzero orbit spectrum $\{1/2,1/2\}$. Nevertheless,
\[
 \boxed{\ L_A^\star=1,\qquad L_B^\star=2.\ }
\]
The first spectrum occupies two real one-dimensional characters; the second occupies an indivisible real rotation plane. A one-dimensional exact action can retain one of the former but none of the latter. Thus regression error and an ordinary spectrum are insufficient summaries of exact-action recovery. The full labeled observations distinguish these examples; this is a separation of summaries, not nonidentifiability of the full data. \Cref{sp:counterexample} gives the calculation.

\subsection{Bounding the entire exact-action class}
Let $\gamma_r$ be the fraction of observation energy captured by the best group-invariant subspace of dimension at most $r$, respecting real irreducible types. Equivalently, it is the type-constrained spectral budget of the orbit Gram matrix, divided by $r$ (\cref{sp:budget,sp:observable}). \Cref{sp:main} proves
\[
 \boxed{\ 2(1-\gamma_r)\le L^\star\le2(1-\gamma_r^2).\ }
\]
The lower bound holds for every member of $\mathcal R_r(G)$, and some member of the same class attains the upper bound or better. For $\gamma_r<1$, the bounds differ by a factor $1+\gamma_r\le2$; both vanish at $\gamma_r=1$. This connects best invariant approximation with fixed-dimensional prediction risk, using classical tools credited in \cref{sp:positioning}. The proof is constructive in an invariant function space; \cref{sec:pending-studies} evaluates its finite-bank action witness and separates construction-bank guarantees from transfer loss.

\paragraph{The cost of the constrained class.}
Because independent least squares contains every exact orthogonal action, $L^\star\ge\eta$. Subtracting $\eta$ from the bounds gives
\begin{equation}
 \max\{0,2(1-\gamma_r)-\eta\}
 \le L^\star-\eta
 \le 2(1-\gamma_r^2)-\eta.
 \label{eq:risk-premium-main}
\end{equation}
This interval quantifies the additional risk of imposing both exact group laws and orthogonality, on the same bank and in the same metric. It is a direct consequence of the risk bound, not an exact measurement of $L^\star$ or a new independent theorem. A lower endpoint above zero rules out an equally accurate exact orthogonal action in that specified class. The identity action $\rho_g=\Id_r$ for every $g$ has risk $L(I)=2(1-e_{\rm triv})$, where $e_{\rm triv}$ is total trivial-type energy. This gives a same-metric reference distinct from the selected budget share $b_{\rm triv}$ (\cref{sp:identity-reference}).

A large $\gamma_r$ guarantees a low-risk exact orthogonal action in the declared metric. It does not distinguish a semantic computation from stable category coding. State-invariant components can contribute to the budget, so total type energy and the selected portion should be distinguished (\cref{sp:type-energy}). These components need not be purely contextual; even nontrivial state structure can reflect category identity. The trivial representation is itself a valid action, so a large trivial contribution limits the semantic interpretation of the score rather than invalidating the action class.

\paragraph{Metric and evidence contract.}
This theoretical risk differs from the historical $D$ and later training-energy-normalized scores. Ridge and finite-step penalized Joint need not be exact orthogonal actions, so their reported scores are not witnesses for this bound. The measured spectra below define finite-bank targets; no independent-context certificate is reported. The conservative certificate in \cref{sp:certificate} requires a sampling contract unmet by the present panel. The $r=1$ example establishes a possible distinction, not its empirical magnitude at $r=16$. For C7/D4, the type correction satisfies $\Delta_{\rm type}/r\le1/16$ at that dimension. In particular, $(1-\eta)-\gamma_r\le1/16$; a large separation in that direction should not be an assumed empirical outcome. The risk-premium interval in \cref{eq:risk-premium-main} is a different quantity. Details are in \cref{sp:type-scale}.

\subsection{Consequences for finite-state recovery}
Three elementary facts explain the empirical controls; their formal statements and proofs are in \cref{app:theory}. First, any $n$ centered affinely independent state vectors in $n-1$ dimensions admit an exact action of every state permutation (\cref{prop:finite}). Seven generic weekday vectors therefore realize a cycle without group-specific semantic learning. A stable state code plus complementary context features can also share one action across contexts. Held-out templates do not exclude that construction, although arbitrary separate contextual realizations need not share a predictor.

Second, composition residuals depend on prediction errors and operator norms (\cref{prop:propagation}). Independent fitting removes group-law penalties from the objective, but does not show that measured consistency adds identification information beyond accurate prediction. The common-design ridge identity in \cref{app:ridge-identity} makes this dependence explicit.

Third, displacements averaged over a complete uniformly weighted state orbit cancel for every permutation (\cref{prop:additive}). The global additive baseline is therefore an identity predictor by construction. Source-conditioned displacements avoid that cancellation and can compose by telescoping centroid differences (\cref{app:additive-composition}). The comparison below evaluates such alternatives with their extra information stated.

The orientation and commutant counterexamples in \cref{app:orientation,app:commutant} further show that equal within-context geometry need not permit common transport. Their restricted constructions delimit interpretation; they do not supply a null distribution for measured LLM errors.

\section{What the Completed Controls Identify}
\label{sec:evidence}
\Cref{tab:main-findings} summarizes the three empirical findings; \cref{tab:identification-summary} retains the detailed endpoint-by-endpoint comparison. Compare losses within each endpoint: historical $D$, compact-coordinate prediction and context-difference transport use different denominators and sampling contracts.

\subsection{Prediction and consistency: two properties, coupled evidence}
Independent ridge attains $\Eeq=0.04830/0.05200$ for $C_7/D_4$, while joint errors are $0.04934/0.05286$. Joint fitting decreases relation defects by $0.01699/0.00728$ at a small direct-fit cost. The maps support approximate predictive action structure without group-law penalties in the independent objective. However, \cref{prop:propagation} explains why this is not automatically an additional semantic discovery beyond accurate prediction. These diagnostics do not establish a comparison matched in prediction error and operator norm.

Generator-only fits predict five excluded nonidentity elements in each group. Their errors are $0.074443/0.056558$, versus directly fitted element errors $0.051124/0.050032$. Correct $D_4$ words outperform matched-length wrong words by $1.097329$ under the recorded contrast. This is meaningful transfer in the operator index; the component states and vocabulary were seen in training. Full structural and generator tables, their historical intervals and depth curves are retained in \cref{app:results}.

\subsection{Invalid registration differs from a coherent alternative}
True-registration homomorphism scores in AFS-B lie below the non-automorphic reference's 2.5\% quantile in all 120 cells per group and metric. Its restricted ensemble excludes the observed registration; no calibrated semantic-null inference is claimed. AFS-C refits Ridge on corrupted training pairs; errors increase against unchanged true held-out targets. Ordinary predictors can degrade similarly, so this establishes sensitivity to supervision rather than semantic specificity.

Study~2 instead consistently supervises an alternative action induced by swapping two state labels. It preserves group laws and fixed-point counts. All-element ridge errors for the specified/coherent correspondence are $0.014897/0.017503$ in $C_7$ and $0.041919/0.041871$ in $D_4$. Each correspondence uses the same underlying source bank and abstract action labels, but its own targets and changed-state evaluation mask. In $D_4$, conjugacy preserves fixed-point counts while potentially changing their identities, so the scored source/action rows need not coincide. This compares the predictability of two registered tasks, rather than two predictors on identical target rows. The single conjugate alternative remains competitive descriptively despite the invalid-registration separation; it does not characterize all semantic alternatives or nonisomorphic representations.

\subsection{Predictor rankings change with information and storage}
\label{sec:characterization}
The nonlinear and category-based predictors below are predictive comparators; their inclusion does not establish exact nonlinear group actions. The theoretical risk comparison in \cref{sec:boundaries} concerns independent linear maps and exact orthogonal actions.

Rank-matched embedding fits are near exact at $q=6$ for weekdays and $q=3$ for corners; internal errors are higher in the two-model lexical panel. \Cref{prop:finite} explains why this can occur without group-specific semantic learning. Equal rank retains the entire lexical span but can discard internal contextual variation, so it does not match retained information or compare all intrinsic structure at the two locations.

In Study~2 at $q=16$, the decoded-centroid predictor improves the $C_7$ mean but wins only 9/24 cells; it is worse on the $D_4$ mean and wins in 2/24. It receives training state labels and the exact action table, beyond source-coordinate/action access. This is an explanatory control with an information advantage. The main table uses the compact extension's equal-action aggregate; the earlier characterization uses a pooled changed-pair aggregate. For D4, differing fixed-point counts make those weights unequal, so the two recorded means need not agree (\cref{app:aggregation-contract}).

The compact nonlinear comparison fits RFF and Nystr\"om residual predictors using all construction target pairs. Neither improves the group mean at the identity-pruned linear storage budget; Nystr\"om does at $2\times$ storage and RFF at $4\times$ in both groups. \Cref{fig:budget-frontier} shows all five declared nominal budgets; exact means and unit-budget comparisons remain in \cref{tab:reproduction-budgets,tab:reproduction-compact}. Storage includes predictor features, normalizers, coefficients and biases, but excludes shared PCA, model weights and runtime memory. The ranking is estimator- and budget-dependent, not universal linear superiority or total-compute efficiency. These nonlinear predictors are shared across contexts, but their prediction scores do not establish invertibility or group laws. Budgets were not selected retrospectively as deployment operating points.

\begin{figure*}[t]
\centering
\includegraphics[width=\textwidth]{figures/budget_frontier.pdf}
\caption{\textbf{Predictor rankings change with storage.} Recorded direct-prediction means for all nonidentity actions on changed states under the specified semantic registration, over 24 fixed model/readout cells per domain. Losses use the Study~2 compact-comparison training-coordinate-energy denominator; lower is better. The horizontal axis is the nominal incremental parameter-storage budget relative to identity-pruned ridge, not runtime memory or statistical capacity. All five declared budgets and all three budgeted predictor families are shown. Nystr\"om first improves the group mean over ridge at $2\times$, and RFF at $4\times$, in both groups. Lines connect recorded aggregates on logarithmic axes; they are not fitted scaling laws or confidence bands. This direct-prediction endpoint is distinct from generator-only composition.}
\label{fig:budget-frontier}
\end{figure*}

\subsection{A transport advantage within the original construction}
Fitting context differences directly yields losses $0.03426/0.11639$ for $C_7/D_4$, versus $0.04714/0.13230$ for applying coordinate-fitted operators to those differences. The coherent action gives $0.03446/0.11717$ under the same direct fitting procedure and registration-specific scoring just described. Within that original bank, the improvement identifies a useful fitting objective; close means for distinct registered tasks establish neither semantic preference nor equivalence. Matched-state context differences cancel state means but retain context main effects, so this endpoint does not isolate state--context interactions. A fitted difference predictor also need not define a group action on the original coordinates. On jointly new $S_3$ contexts/entities, direct transport loss is $0.016290$ versus identity $0.016444$, a small absolute improvement below the inherited $0.01$ practical margin. Complete normalization and information conditions are in \cref{app:characterization,app:reproduction}. The new authored bank in \cref{tab:followup-context} tests this advantage under matched validation selection and reveals a domain-dependent reversal.

\paragraph{Behavioral scope.}
Historical activation edits supply limited behavioral evidence: $C_7$ target logit improves over identity by 1.129, but the mean effect is smaller than that of the source-conditioned additive control. Correct $D_4$ ordering has lower final-logit distance than the matched wrong ordering. These are token/logit outcomes, not a certified native algorithm. The additive control uses full-hidden displacements and source identity, whereas the operator edits a PCA subspace; the historical mean-overwrite hook lacks a matched zero-delta sham. The separate native program branch failed clean competence and its confirmation remains blocked. Full endpoints and limits, including the unchanged failure records, are in \cref{app:behavior,app:reproduction}.

\input{sections/planned_studies.tex}

\section{Implications for Evaluating Recovered Actions}
\label{sec:reading-protocol}
The findings motivate four complementary questions, conditional on states, contexts, coordinates, estimator and outcome. They organize interpretation rather than define a calibrated certification test or four statistically independent sources of evidence.

\paragraph{1. Predictive recovery.}
State what was held out: contexts, action elements, state identities or domains. Disclose comparator supervision and storage. Our generator result demonstrates why the distinction matters: transfer to an excluded action element can still use a completely familiar state vocabulary.

\paragraph{2. Action consistency.}
Specify the group-law residual, its metric and whether fitting encouraged it. The propagation identity relates consistency to prediction errors and operator norms. To claim additional identification information, compare predictors under conditions that account for those quantities; independent objectives alone do not establish independent evidence.

\paragraph{3. Preferential semantic alignment.}
Compare coherent alternatives as well as invalid registrations, making information, capacity, scoring rows, target definitions and predictive difficulty explicit. Our comparisons show why these controls answer different questions. Preference over selected alternatives would narrow the interpretation, rather than identify a unique semantic mechanism.

\paragraph{4. Native causal use.}
For a claim about successful native computation, establish clean-task competence and intervention fidelity, then distinguish sufficiency from necessity for the proposed mechanism. The historical edits illustrate the gap between moving an output score and identifying the computation used by the unedited model.

For a new application, specify the intended claim and competing explanation before choosing a diagnostic. The useful outcome is the interpretation changed by the comparison, including unfavorable results, rather than the number of positive scores obtained.

\section{Discussion and Scope}
The results connect three distinct targets. Predictors recover useful transformations; the spectral bound quantifies compatibility with an exact orthogonal action; exhaustive image comparisons ask how special the specified registration is within a finite comparator class. The finite-bank interval constrains the optimum over the entire exact orthogonal-action class in its declared metric. Frozen-witness losses measure one constructor's transfer performance, while registration margins compare spectral budgets across the fixed image class. Their answers need not agree. In particular, the constructed action is feasible but less accurate than OLS in every held-out condition, and a high spectral budget is often shared by other registrations. The observed real-type correction is negligible in 47 of 48 conditions at rank 16, so the analytical separation is not a large universal empirical effect in this panel.

Context transport supplies a further practical distinction. An objective that improves both domain means within the original construction improves only the weekday mean under the new construction, with square-corner results reversing. This prevents treating dedicated difference fitting as a universally better intervention or transport method. It also motivates reporting domain and model variation rather than only the nearly cancelling pooled contrast.

The evidence is restricted to selected base models, readouts and small seen-state vocabularies. New contexts are assistant-authored paired frames with operator review; both the outer narratives and limited factual wording differ from the old bank. They are not independently sampled human prose, and at a causal decoder's state-token readout only text preceding or within that span can affect the representation. The retained S3 extension includes new entities and contexts but yields only small gains over identity. No native causal claim is made: clean-competence and other original scientific stops remain in the record. The finite-bank theorem and fixed-mechanism illustrations do not calibrate a population-level identification procedure. These limits delimit the contribution without erasing its positive content: a computable action-class target, evaluated feasible witnesses, and complete, interpretable small-group comparisons.

\section{Conclusion}
Predictive recovery, exact-action compatibility and semantic specificity are different properties. Our bound makes the second measurable on a complete finite bank, while the witness and exhaustive image comparisons show why it does not settle the third. The new-context study additionally demonstrates that a transport objective's advantage can depend on the language construction and domain. Together, these results support evaluating recovered actions through explicit action classes, metrics and alternatives, rather than treating low prediction and consistency errors as a single identification claim.

\FloatBarrier
\label{sec:main-end}

\section*{Impact Statement}
This work develops diagnostics for interpreting language-model representations. Overinterpreting probes can mislead users about model capabilities or mechanisms. These controlled tasks do not assess demographic fairness or establish safety guarantees. We identify no other specific societal consequences requiring discussion.

\bibliography{references}
\bibliographystyle{tmlr}
\appendix
\input{sections/theory.tex}
\input{sections/spectral_risk.tex}
\input{sections/methods.tex}
\input{sections/results.tex}
\input{sections/characterization.tex}
\input{sections/reproduction.tex}
\input{sections/prospective.tex}
\input{sections/related_work.tex}
\input{sections/planned_protocols.tex}
\end{document}
